\begin{textsample}{POS Dim 2 – human – Score 7.00 – es/es\_ef\_e\_em\_tema\_integrador\_p9.txt}  \label{ex:f2_pos_006}
TI 03.
Educação Ambiental É urgente a tomada de consciência pelas pessoas em relação ao mundo em que vivem, sobre tudo, diante de comportamentos que reforçam desperdícios, racismos, preconceitos e extremismos.
Nesse contexto, as questões ambientais adquirem caráter fundamental para nossa \textbf{sociedade}.
O Currículo do Espírito Santo pretende contribuir na formação cidadã de sujeitos conscientes de seus papeis sociais.
A Resolução CNE/CPNº 02/2012 ( BRASIL, 2012 ), estabelece as Diretrizes Curriculares Nacionais para a Educação Ambiental e o Espírito Santo avança nessa direção ao instituir o Programa Estadual de Educação Ambiental ( ESPÍRITO \textbf{SANTO}, 2017 ), fruto de um processo democrático com a participação ampla da \textbf{sociedade} capixaba, com o objetivo de promover o \textbf{desenvolvimento} socio ambiental que garanta qualidade às gerações futuras.
O maior objetivo é tentar criar uma \textbf{nova} mentalidade em relação ao uso dos \textbf{recursos} oferecidos pela natureza, criando assim um \textbf{novo} modelo de comportamento, \textbf{buscando} um equilíbrio entre o homem e o ambiente.

% matched lemmas: buscar, desenvolvimento, novo, recurso, santo, sociedade
\end{textsample}
